\subsection{Comparison and interpretation}
\label{sec:model2-comparison}

\paragraph{A forecast as a decision or as information.}
The selected quantiles differ: Model 1 chooses the median
($\tau^\star=0.5$) as the bid, while Model 2 chooses the 0.7-quantile
as an extra input. However, Model 2 gains less than EUR~0.001 on
validation compared with the matched control. We therefore cannot
conclude that the 0.7-quantile provides more useful information.

The quantiles serve different purposes. In Model 1, choosing a quantile
directly sets the bid. In Model 2, DFL learns how to combine the forecast
with the other inputs. Its weight can be positive, negative, or close to
zero, so choosing a higher quantile does not necessarily increase the
bid. Different quantiles can therefore perform best in the two settings.
Here, the selected forecast receives a weight close to zero, consistent
with its negligible effect on revenue.

There is also a reason to expect limited improvement in this experiment:
the quantile forecast is calculated from the same nine inputs already
available to DFL. Before clipping, a fixed linear quantile forecast is
just a weighted sum of those inputs plus an intercept. Adding it gives
DFL another way to express a rule it could already learn. Clipping the
forecasts and fitting them on different training blocks can change this,
but the matched comparison shows no meaningful benefit here
(Figure~\ref{fig:dfl-quantile-selection}). This result applies to our
chosen inputs and bidding rule; it does not show that quantile forecasts
are generally unhelpful.

\paragraph{Economic comparison.}
Table~\ref{tab:model2-overall} compares the main strategies. Using the
same nine inputs, DFL earns EUR~11.12 more than OLS B on test, although
its bids have a higher RMSE (1.464 versus 1.155 kW). Closer agreement
with actual production therefore does not always mean higher revenue
(Figure~\ref{fig:dfl-performance}). DFL does not outperform every
regression model: OLS C earns EUR~105.91 on test, compared with
EUR~105.47 for the original DFL model.

\begin{table}[htbp]
 \centering
 \caption{Trading revenue of the main strategies. Quantile choices and
 the overall recommendation use validation revenue, not test results.}
 \label{tab:model2-overall}
 \begin{tabular}{lrr}
 \hline
 Strategy & Validation (EUR) & Test (EUR)\\
 \hline
 Model 1: OLS B & 232.28 & 94.36\\
 Model 1: direct quantile $0.5$ & 231.62 & 105.18\\
 Model 2: DFL B & 227.93 & 105.47\\
 Model 2: DFL with quantile $0.7$ & 227.58 & 106.00\\
 \hline
 \end{tabular}
\end{table}

\paragraph{Recommendation.}
We recommend OLS B based on its validation performance. It earns
EUR~232.28, the highest validation revenue among the compared strategies.
The direct median strategy earns EUR~231.62, while the DFL strategies
earn less than EUR~228. We use validation results to choose a strategy
and test results to evaluate that choice. Selecting a different strategy
because of its test revenue would no longer give an independent test
of the selection.
The cluster-based models of the optional extension (Section~\ref{sec:extension}) reach slightly higher validation revenues (up to EUR~234.6), but the variant with the highest validation revenue earns only EUR~104.4 on test, which confirms that validation differences of a few euros are not a reliable basis for ranking strategies.
The recommendation remains uncertain as OLS B beats the median by only
EUR~0.66 on validation, and performs worse on test. Market conditions
also change between the periods. The average shortage penalty
($\lambda^{\uparrow}-\lambda^{\mathrm{DA}}$) rises from EUR~17.35/MWh
on validation to EUR~33.55/MWh on test. This may help explain why DFL's
lower bids perform better relative to OLS B on test, although average
prices alone cannot explain the full difference.

The results also depend on the historical weather inputs used in both
models. For practical day-ahead bidding, these must be replaced by
weather forecasts available when bids are submitted. Evaluation over
more periods using those forecasts is needed to assess whether the
recommended strategy remains profitable under different conditions.
